%% Secciones comunes a todos los manuales, edición española.
%% Se incluyen con \input{common-es}. Lo único que cambia entre applets es
%% \appletfolder, que es también el nombre de la carpeta y de la ruta en el sitio.

\section{Cómo abrirlo}

No hay que instalar nada. La aplicación entera es un solo archivo
\code{index.html}: no descarga librerías, no llama a ningún servidor y no
guarda nada fuera del navegador.

\subsection{Las tres maneras}

\begin{description}
\item[Doble clic sobre el archivo.] Es lo más rápido. Busque
  \code{\appletfolder/index.html} y ábralo. El navegador lo muestra desde el disco,
  con la barra de direcciones empezando por \code{file://}. Funciona sin
  conexión desde el primer momento.
\item[Desde la web.] Si alguien ya lo publicó, la dirección termina en
  \code{/\appletfolder/}. Una vez cargada la página, puede desconectarse: no
  necesita la red para nada más.
\item[Servido desde una carpeta.] Para una clase, poner la carpeta en
  cualquier servidor estático basta. Con Python instalado:
  \begin{quote}\code{python3 -m http.server 8000}\end{quote}
  y luego \code{http://localhost:8000/\appletfolder/} en el navegador.
\end{description}

\begin{amnote}
Guardar la página con \key{Ctrl}\,+\,\key{S} (o \key{Cmd}\,+\,\key{S} en
Mac) deja una copia propia que sigue funcionando sin conexión y que se puede
llevar en una memoria USB o repartir por correo.
\end{amnote}

\subsection{Idioma y tema}

Arriba a la derecha hay dos pares de botones.

\begin{ctrltable}{Control & Qué hace}
ES / EN & Cambia el idioma de toda la interfaz al vuelo. Nada se recalcula:
          puede cambiarlo en mitad de una explicación sin perder el estado. \\
Oscuro / Claro & Cambia el tema. El oscuro es el diseño original; el claro
          existe para proyectores y para imprimir. \\
\end{ctrltable}

Las dos elecciones se recuerdan en el navegador, así que la próxima vez se abre
como la dejó.

\subsection{Qué navegador}

Cualquiera que sea razonablemente reciente: Chrome, Edge, Firefox o Safari, en
ordenador, tableta o teléfono. En pantallas estrechas los paneles se apilan uno
sobre otro en lugar de ponerse al lado. No hace falta cuenta, ni permisos, ni
extensiones.
